%last updated in April 2002 by Antje Endemann 

 
\documentclass[runningheads]{llncs} 
%In order to omit page numbers and running heads 
%please use the following line instead of the first command line: 
%\documentclass{llncs}. 
%Furthermore change the line \pagestyle{headings} to 
%\pagestyle{empty}. 
       
\usepackage{pstricks,pst-node}

\begin{document}
 
 
   
\newlength{\TTL}
\newlength{\TMPL}
\newcommand*{\Let}[2]{\ensuremath{#1\Leftarrow#2}}
\newcommand*{\AlgoStep}[1]{Step {#1}}
\newcounter{AlgoStepCnt}
\newcounter{algo}
\makeatletter
\newenvironment{BAlgo}[2][1]{%
%  \newcommand*{\Let}[2]{##1\Leftarrow##2}
  \newcommand*{\For}[2]{\textbf{for} $##1$ \textbf{to} $##2$}
  \newcommand*{\If}{\textbf{if }}
  \newcommand*{\Then}{\textbf{then }}
  \newcommand*{\Else}{\textbf{else}\xspace}
  \newcommand*{\End}{\} }
  \newcommand*{\Do}{\textbf{do} \{ }
  \newcommand*{\Return}{\textbf{return}\xspace}
  \newcommand*{\Until}[1]{\} \textbf{until} ##1}
  \newcommand*{\While}[1]{\} \textbf{while} ##1}
  \newcommand*{\WhileDo}[1]{\textbf{while} ##1 \Do}
  \newcommand*{\T}[1]{\rule{3ex}{0pt}\setlength{\TMPL}{\TTL}\addtolength{\TMPL}{
-3ex}\parbox[t]{\TMPL}{\setlength{\baselineskip}{0pt}{##1}}}
  \newcommand*{\TT}[1]{\rule{6ex}{0pt}\setlength{\TMPL}{\TTL}\addtolength{\TMPL}
{-6ex}\parbox[t]{\TMPL}{\setlength{\baselineskip}{0pt}##1}}
  \newcommand*{\TTT}[1]{\rule{9ex}{0pt}\setlength{\TMPL}{\TTL}\addtolength{\TMPL
}{-9ex}\parbox[t]{\TMPL}{\setlength{\baselineskip}{0pt}##1}}
  \begin{figure}[htb]\refstepcounter{algo}%
   % \setlength{\textwidth}{8.6cm}%
  \begin{minipage}{\textwidth}%
      \centerline{Algorithm~\thealgo: \textbf{#2}}\vspace{-1ex}%
       \centerline{\rule{0.93\textwidth}{0.5pt}}%
      \begin{list}{}{\setlength{\itemsep}{0pt}\setlength{\parsep}{3pt}%
          \setlength{\topsep}{0pt}\setlength{\parskip}{0pt}\setlength{\partopsep
}{0pt}%
          \setlength{\baselineskip}{0pt}%
          \setlength{\leftmargin}{2em}%
          \setcounter{AlgoStepCnt}{#1}%
          \setlength{\rightmargin}{1em}%
          \setlength{\TTL}{\textwidth}%
          \addtolength{\TTL}{-3em}                 
          \renewcommand{\makelabel}{\hfill\textit{\footnotesize{\arabic{AlgoStepCnt}}}\stepcounter{AlgoStepCnt}~}%
          } }%
      {\end{list}
 
      \centerline{\rule[1ex]{0.93\textwidth}{0.5pt}}
      \end{minipage}\end{figure}}




\mainmatter 
 
\title{An analysis of the performance of the Mixture of Trees Factorized Distribution Algorithm when
  priors and adaptive learning are used} 
 
\titlerunning{Lecture Notes in Computer Science} 
 
\author{Roberto Santana \inst{1}} 
 
\authorrunning{Roberto Santana} 
 
\institute{Institute of Cybernetics, Mathematics, and Physics (ICIMAF),\\  
Calle 15, entre C y D, Vedado, C-Habana, Cuba\\ 
\email{\{Rsantana\}@cidet.icmf.inf.cu}\\ 
} 
 
\maketitle 
 
\begin{abstract} 
  This paper analyzes the behavior of the Mixture of Trees Factorized
  Distribution Algorithm (MT-FDA) when priors are incorporated. It is shown
  that the addition
  of priors provokes a mutation like effect during the search. Adaptive priors that relate the rate of
  mutation to the quality of the search are also introduced. Additionally, the
  learning step of the MT-FDA is changed to avoid the overfitting of data. The
  results of the experiments show that our proposals improve the trade off
  between exploration and exploitation displayed by the MT-FDA. 

\end{abstract} 
 
 
\section{Introduction} 
 


Population Based Search Methods that use Selection  (PBSMS) are non deterministic
heuristic search strategies, commonly used as optimization methods. Their main
characteristics are: They use a population of points instead  of a single
point to conduct the search. In every iteration (usually called generation) a
subset of points is  selected, and  by applying some operators a new population is created. In  this way the algorithm iterates (evolves) until one stop condition is  satisfied.  
  
    Genetic Algorithms \cite{Holland:1975} are members of a subclass of
    PBSMS where recombination and mutation operators are applied to the selected  set
    of individuals to obtain the new population. Another class of PBSMS  comprises
    those algorithms    characterized by the use of probabilistic modeling of the information  contained
    in the selected set, instead of the genetic operators. These  algorithms are the
    focus of this paper. 

 In this paper  an
   Estimation Distribution Algorithm is any evolutionary algorithm
    that uses the estimation of probability distributions to improve
    the search. Although no all the
   proposals fit well in the EDAs scheme this conceptual framework has been taken before
   \cite{Larranhaga_et_Lozano:2001}
    as a model to study the algorithms  under analysis. A special subclass of EDAs will group the
    algorithms that use factorizations of the probability
    distribution. This subclass of algorithms will be called Factorized
    Distribution Algorithms (FDAs)\footnote{In the literature the term FDA
    is usually given to a Factorized Distribution
    Algorithm that uses the same factorization along the evolution}.


In this paper we concentrate on the  Mixture of Trees FDA (MT-FDA) that was
 introduced in \cite{Santana_et_al:2001b}. The
 algorithm uses as its probabilistic model a mixture of trees. Our main contribution in \cite{Santana_et_al:2001b} was to show that a FDA that uses
as its probabilistic model a  mixture of trees can perform better than other simpler
FDAs and be competitive, and some times superior to FDAs that use more complex
probabilistic models. In the present study we show how the results of the MT-FDA can
 be improved by adding a mutation like effect using priors, and modifying
 the algorithm used for learning the probabilistic model.


 Along the paper we will concerned with the
 maximization of a function $f: B^n \rightarrow R$,  and  $X \in  B^n$  is a set
 of random binary variables. We will use $x_i$ ($x_i \in \{0,1\}$) to denote a
value of $X_i$, the i-th component of $X$.  $P=\{x^1,\ldots x^N\}$ will
 denote a set of vectors where $N$ is the number of binary
 vectors in the set.
 
 The paper's outline is as follows: In the next section we will present the mixture of trees as a probabilistic
model and the MT-FDA. In section~\ref{EXPEXP} we discuss the relationship that
exits in EDAs between the goals of
exploration and exploitation and probabilistic
modeling. Section~\ref{PRIORS} describes how priors can be inserted in the
MT-FDA to obtain a mutation like effect. Adaptive priors are defined for the
first time in EDAs. Section~\ref{ALEARNING} presents a modification to the learning
algorithm based on halting the learning procedure to avoid overfitting. Experimental results on the application of the proposed
modifications  are shown in section~\ref{EXPERIMENTS}. The conclusions of the paper and  future research trends
are presented in section~\ref{CONCLUSIONS}.

\section{Trees and Mixture of Trees Models} 

 Modeling by finite
mixture of distributions \cite{Everitt_and_Hand:1981} concerns modeling a statistical
distribution by a mixture (or weighted sum) of other distributions. The mixture of
trees  was introduced as a probabilistic model in \cite{Meila:1999} where its
usefulness as a density estimator algorithm and a classification tool was
demonstrated on a set of problems. Mixture of trees, and in general mixture distributions have a number of 
distinctive attributes that make them particularly appealing for their use 
in the framework of FDAs. Maybe the most important is the possibility of 
representing, condensed in just one model, different patterns of 
interactions among the variables of the problem. In Bayesian Networks (BN) the change in one 
variable's value can determine changes only in the parameters of other 
variables, no in their structural relation (edges or arcs in the graphical 
representation). In mixture of finite distributions the structure of 
dependencies among a set of variables can change depending on the values of 
the choice variable they depend on.

 
 Now the probabilistic models are 
formally introduced. We will utilize the same notation used in \cite{Meila:1999}. 
  A probability distribution $T$ that is conformal with a tree is defined as: 
 
\begin{equation} 
   T(x)=\prod_{v\in V} T_{v|pa(v)}( x_v|x_{pa(v)}) \;.  \label{T-Tree} 
\end{equation} 
 
The distribution $T$ itself will be called a tree when no confusion is 
possible. The graph $(V,E)$ represents the structure of the distribution $T$. A mixture of trees is defined to be a distribution of the form:



\begin{equation}
   Q(x)=\sum\limits_{k=1}^m\lambda _kT^k(x) \;  .  \label{Q-Mixture}
\end{equation}

with $\lambda _k\geq 0$, $k=1,\ldots,m$, $\sum_{k=1}^m\lambda _k=1$.


The tree distributions are the mixture components, and the $\lambda _k$ are called
mixture coefficients. A mixture of trees can be viewed as containing an unobserved
choice variable $z$, which takes values $k\in $ $\{1,\ldots,m\}$ with probability
$\lambda _k$. Conditioned on the value of $z$ the distribution of the visible
variables $V$ is a tree. The $m$ trees may have different structures and different parameters.

 \begin{BAlgo}{MT-FDA}
 \label{MTFDAALG}
 \item Set $t\Leftarrow 0$. Generate $N\gg 0$ points randomly. 
 \item \Do
 \item \T {Select a set $S$ of $k$ $\leq N$ points according to a selection method.} 
 \item \T {Calculate a mixture of trees $Q$ that approximates the distribution
     of points in $S$.}
 \item \T {Generate new points sampling from $Q$.}  
 \item \T {$t \Leftarrow t+1$} 
 \item  \Until{Termination criteria are met.} 
\end{BAlgo} 
 
 
Above, the pseudo-code of the MT-FDA is presented. The first population of 
the algorithm is generated randomly. From the current population, a subset 
of points is selected. A mixture of trees $Q$ that fits the selected set is 
found using the Iterated Estimation Maximization (IEM) algorithm \cite{Meila:1999}. New points are generated by sampling from
$Q$. Different replacing strategies can be used 
to combine points from the current population and new generated points in the next
population.  When proportional or Boltzmann
selection are used the step 3 of algorithm~\ref{MTFDAALG} is not needed and in
the step 4 of the
algorithm  the model is calculated using the probabilities of selection
associated to each point in the population by the selection method. This can be done because graphical
  models can be learned too from a joint probability distribution.


\section{Exploration and Exploitation in EDAs} \label{EXPEXP}
 

 When generating a new population the two traditional goals of an efficient search, exploitation 
and exploration of the search space, have to be accomplished. In the case of FDAs
the achievement of these goals will depend on the power of expression of the probabilistic model, the way it has been learned and the method used to sample
points from it. In principle, a good probabilistic model has to be able to exploit
 the good points that have been already identified, in our case represented by the
 set of selected points and possibly some elite points, but also able  to
 explore new areas of the search space where eventually are located better
 points. 

 In the model learning phase of Bayesian FDAs \cite{Pelikan_et_al:1999,Larranhaga_et_Lozano:2001} the quality
 of the model is usually evaluated  using different measures of the likehood of the
 data of the selected points. It is clear that a probabilistic model with a
 maximal likehood of the data can be useless to explore new areas of the search
 space. While an exact learning
of the probabilistic model from the data can benefit exploitation, if learning
is pursued beyond certain limits the phenomenon of overfitting arises. An overfitted  model  can less likely generate, during the sampling step, points that belong 
to unexplored areas of the search space.  Different alternatives have been
 proposed to cope with this problem in probabilistic modeling. In the context of FDAs
 we identify the following possible solutions to avoid overfitting:

  \begin{itemize}
  \item To allow the learning algorithm to learn a very exact (possibly overfitted)
    model and modify this model after.
  \item To stop the learning process when a predefined quality in the approximation
     of the data has been achieved. 
  \item To do the search of the probabilistic model in a constrained space of
     (simple) models.
    
  
  \end{itemize}


 The first idea has been implemented for the Univariate Marginal Distribution
 Algorithm (UMDA) \cite{Muhlenbein_and_Paas:1996}, the Factorized Distribution
 algorithm with a fixed model ($FDA^*$) \cite{Muhlenbein_et_al:1999}
  and Learning FDA (LFDA) \cite{Muehlenbein_et_Mahnig:2001a} by means of considering
 Bayesian priors during the learning of the model
 \cite{Mahnig_and_Muehlenbein:2001a}. In the implementation, the learned
 model's parameters are changed after the structure of the model has been learned. The effect of this type of priors is similar to
 the effect of mutation in Genetic Algorithms.  The authors show that, also for FDAs,
 mutation increases in many cases the performance of the algorithm. This can be
 explained by the important role played by mutation in avoiding premature convergence
 and allowing the exploration of new
 regions during the search. In our first implementation of the MT-FDA we did not consider the use of
 priors. One of the improvements incorporated to the algorithm in  this paper is the
 use of priors.
 
Although the use of priors can in general improve  the results of the search, it
has  been shown in \cite{Mahnig_and_Muehlenbein:2001a} that it can decrease
the quality of the results for certain functions. This fact makes convenient to
investigate the second alternative mentioned before, i.e. to stop the learning process when a predefined quality in the approximation
     of the data has been achieved. In section~\ref{ALEARNING} we analyze this issue.

\section{Use of Priors}\label{PRIORS}

 In
\cite{Mahnig_and_Muehlenbein:2001a} it is explained how to introduce mutation into
FDAs and how to choose the mutation rate based on a theoretically derived
result.  In
the paper it is shown that mutation increases in many cases the performance of the
algorithms and decreases the dependence on the correct choice of the population
size.  The authors state that:

Let $\tau$ denote the parameter for truncation selection, $I_{\tau}$ is the
strength of selection that depends on $\tau$ \cite{Muhlenbein:1997}, and $M$ is
the size of selected set. When $r$ is the prior for a single binary variable, the prior $r'$ for a factor
$p(x_1,\ldots, x_k)$, and the prior $r^*$ for a factor $p(x_k|x_1,...,x_{k-1})$ should be

\begin{equation}
   r'=r^*=2^{-(k-1)}\cdot r \;. 
 \label{GENPRIOR} 
\end{equation}
 
Using $r'_i=2^{-k_i-1}r$ with $r=\frac{I_{\tau}M}{n}$ is a reasonable choice for the
  Bayesian prior for truncation selection
  \cite{Mahnig_and_Muehlenbein:2001a}. 


We
  use the proposed prior in our experiments adapting it to the case of the MT-FDA
  where the marginals used can have only order one or two. Additionally, we have introduced
  an adaptive prior that can be different for each tree. This prior increases the probability of the appearance of events with
  zero probability proportionally to the approximation of the data, and inversely 
proportionally to the coefficient of the tree (i.e. its weight in the mixture
  approximation). 

 Let us introduce some new notation.

We define $D^c$ as the set formed by exactly one copy of all the different vectors in
$D=\{x_1,x_2,\ldots,x_M \}$. This means $x \in  D^c \Longrightarrow x \in  D$ and $x_i,x_j \in D^c
\Longrightarrow x_i \neq x_j$. We denote as
${\tilde{P}}^k(D^c)$ to the sum of the probabilities assigned by the tree $T^k$ to all the
points in $D^c$. Respectively ${\tilde{P}}(D^c)$ is the sum of the probabilities assigned
by the mixture to all the points in  $D^c$.


\begin{equation}
 {\tilde{P}}^k(x_i) = \lambda_k T^k(x_i) \;. 
\end{equation}

\begin{equation}
 {\tilde{P}}(D^c) = \sum_{x_i \in D^c} \sum_{k=1}^{m} {\tilde{P}}^k(x_i) \;. 
 \label{PMEASURE}
\end{equation}

 $\tilde{P}$ is not a probability in $D^c$ (e.g. $\tilde{P}(D^c) \neq 1$). Only
 when all the points of the search space are in $D^c$, $\tilde{P}(D^c) = 1$. The prior $r^k$ for the $k$ tree is  defined as follows.

\begin{equation}
   r^k =  \frac{\tilde{P}(D^c)M}{{\lambda}^kn}
 \label{ADAPTPRIOR} 
\end{equation}

 Finally, we substitute $r$ by $r^k$ in~\ref{GENPRIOR}. This choice of the priors is related to the adaptive mutation schedules used in Genetic
 Algorithms.  


\section{Adaptive Learning} \label{ALEARNING}

We address now the problem of how precise has to be the approximation given by the
mixture of trees in order to avoid overfitting. The IEM algorithm, used in
\cite{Meila:1999} to learn the mixtures, improves the likehood of the data in
each iteration by modifying the structure and parameters of  the mixture. In the
experiments done in \cite{Santana_et_al:2001} we have confirmed that stopping the learning process just when 
no improvement is detected in the likelihood can provoke a premature 
convergence of the MT-FDA because the learned model overfits the data.

Now we present a solution to the problem of determining the extent of
learning. As a measure for assessing the
quality of the learned model in every step of the mixture of trees
learning algorithm we will use $\tilde{P}$, previously defined in equation~\ref{PMEASURE}. $\tilde{P}$ gives an idea of which is the probability
given by the model to points that are not in the data set $D$, $\tilde{P}(x \in X, x \notin D) = 1-\tilde{P}(D^c)$.  So, in principle we can run the learning algorithm while the
probability given by the model to points that are in $D$ does not exceed a given
parameter $\mu$. $\mu$ defines a measure of overfitting, or in  terms of a population
based search method, a measure of exploitation

  When the model is a bad
 approximation of the data, or the data points are a very small sample of the state
 space $\tilde{P}(D^c) \approx 0$. When the model completely overfits the data points
 $\tilde{P^c}=1$.  To summarize, we present the modification done to the IEM learning algorithm: The
 calculation of $\mu$ is incorporated and the condition $\tilde{P}(D) \geq \mu$ is
 added to the termination criteria.

\section{Experiments} \label{EXPERIMENTS}

In our experiments  we will evaluate the impact of all the introduced proposals  in the performance of the MT-FDA. We will focus on the analysis of the
algorithm's behavior when the complexity of the model is constrained  using
$\mu$ as a parameter. Additionally we study the influence of
the number of trees and the use of priors. The test functions used were:
$f_{deceptive3}$, $f_{deceptive4}$, $F_{IsoP}$, $f_{3deceptive}$, and
$Isotorus$. Deceptive functions have served to study the performance of GAs in
the presence of deception. The $F_{IsoP}$ and $Isotorus$ functions
respectively allow the
study of problems with a chainlike and grid based structure. The interested
reader is referred to  \cite{Larranhaga_et_Lozano:2001} for an account of the performance of other FDAs for these
functions. All the functions are defined below and in every case
$u=\sum_{i=1}^{n} x_i$.


Function $f_{dec}^3$: 
 
\begin{equation} 
f_{dec}^3=\left\{  
\begin{tabular}{ccc} 
$0.9$ & $for$ & $u=0$ \\  
$0.8$ & $for$ & $u=1$ \\  
$0.0$ & $for$ & $u=2$ \\  
$1.0$ & $for$ & $u=3$ 
\end{tabular} 
\right. \;.  \label{GOLDBERGDEC3} 
\end{equation} 
 
Function $f_{3deceptive}$: 

\begin{equation}  
 f_{3deceptive}(x)=\sum\limits_{i=1}^{i=\frac n3}f_{dec}^3(x_{3i-2},x_{3i-1},x_{3i}) \;. 
 \label{DEC3}
\end{equation} 

Function general deceptive of order $k$, $fdecK$: 
 
\begin{equation} 
fdecK(x)=\left\{  
\begin{tabular}{ccc} 
$k-1$ & $for$ & $u=0$ \\  
$k-2$ &       & $u=1$ \\  
      & $\cdot \cdot \cdot $ &  \\  
$k-i-1$ &  & $u=i$ \\  
& $\cdot \cdot \cdot $ &  \\  
$k\cdot n$ & $for$ & $u=k$ \\ 
\end{tabular} 
\right. \; .   \label{GDECK} 
\end{equation}
 
Function $f_{deceptiveK}$: 

\begin{equation}  
 f_{deceptivek}(x)=\sum\limits_{i=1}^{i=\frac{n}{k} }fdecK(x_{ki-k+1},\ldots,x_{ki}) \; . 
 \label{FDECK}
\end{equation} 


Function $IsoPeak$: 
 
\begin{eqnarray}
&&
\begin{tabular}{|c|c|c|c|c|} 
\hline 
$x$ & $00$ & $01$ & $10$ & $11$  \\ \hline 
$IsoC_1$ & $m$ & $0$ & $0$ & $m-1$  \\ \hline 
$IsoC_2$ & $0$ & $0$ & $0$ & $m$  \\ \hline 
\end{tabular} 
\label{ISOPEAK} \\
  F_{IsoP}(x) &=& \sum_{i=1}^{m-1} IsoC_1(x_{2i-1},x_{x2i}) +
  IsoC_2(x_{2m-1},x_{x2m}) \;. 
\nonumber 
\end{eqnarray} 
 
with $n=m+1$.   



Function $Isotorus$: 
 
\begin{eqnarray} 
&&  
\begin{tabular}{|c|c|c|c|c|c|c|} 
\hline 
$u$ & $0$ & $1$ & $2$ & $3$ & $4$ & $5$ \\ \hline 
$IsoT_1$ & $m$ & $0$ & $0$ & $0$ & $0$ & $m-1$ \\ \hline 
$IsoT_2$ & $0$ & $0$ & $0$ & $0$ & $0$ & $m^2$ \\ \hline 
\end{tabular} 
\label{ISOTORUS} \\ 
F_{IsoTorus}(x) &=&IsoT_1(x_{1-m+n},x_{1-m+n},x_1,x_2,x_{1+m})+  \nonumber \\ 
&&\sum\nolimits_{i=2}^nIsoT_2(x_{up},x_{left},x_i,x_{right},x_{down})  \;. 
\nonumber 
\end{eqnarray} 
 
where $x_{up}$, etc., are defined as the appropriate neighbors, wrapping 
around.

In the experiments the number of variables is fixed to $36$, as well
as the selection algorithm (Truncation selection with truncation parameter
$0.15$). 



\subsection{Numerical Results}



Table~\ref{MT6FUNCS1} presents the results of the MT-FDA with
$6$  trees when the parameter  $\mu$ is changed and two different types of
priors are used: the recommended and adaptive priors. The table shows the
number of trials (of 100) where the optimum was found  and the average number
of evaluations. For the deceptive functions it is
evident the influence of parameter $\mu$ in the behavior of the algorithm. For
these functions a low value of $\mu$ gives better results than when a higher one is
used.  For the Iso
functions the same behavior can not be clearly appreciated, particularly for
function $F_{IsoTorus}$. 

For deceptive functions the priors increase the number
of times the optimum is found, although in the case
of the adaptive prior there is an interaction with the parameter $\mu$, when  $\mu$
is high the successful rate diminishes.  The
increment in the number of successful trials reached by using priors is
achieved at the cost of an increment in the number of function
evaluations. In the case of the Iso functions there is
not a significative difference when priors are added. Even more, in the case
of function $IsoPeak$ results deteriorate.


In order to create an algorithm of practical use we would like to reduce as much as
possible the number of parameters or to find rules of thumb that can be used to assign
their values.  We hypothesize that a good choice for parameter $\mu$ is around
$0.2$. Notice that the case when $\mu=1$ is equivalent to run the IEM until no further
improvement in the likehood is achieved (as we did in previous implementations).  In the first generations, when the data is
very diverse a good approximation is very hard to reach. The same happens when a high
prior is used. In these cases the learning algorithm will stop when no further
improvement in the likehood  is achieved, or the improvement is under a given
threshold (we use $0.005$ as the threshold value).

\begin{table}[htbp ]
  \begin{center}
  \caption{Results of the MT-FDA  with 6 components in the optimization of different   functions,  $n=36$, $T=0.15$}
   $\begin{array}{|r|r|r|r|r|r|r|r|}
\hline
  Function & Popsize &     \multicolumn{2}{|c|}{No Prior}  &  \multicolumn{2}{|c|}{Best Prior} &
      \multicolumn{2}{|c|}{Adap. Prior}  \\ \hline        
        &   \mu      & succ.   & eval. & succ.  & eval. & succ.  & eval. \\ \hline \hline
 f_{deceptive3} & 0.2 & 80       & 4807 & 100 & 5271 & 100   & 5579           \\ \hline
                & 0.4 & 90       & 4847 &  99 & 5310 & 100   & 5963           \\ \hline
                & 0.6 & 82       & 4979 &  99 & 5296 & 100   & 5886           \\ \hline
                & 0.8 & 67       & 4967 & 100 & 5683 & 99   &  7831           \\ \hline
                & 1.0 & 41       & 4928 &  98 & 5643 & 65   &  7034           \\  \hline \hline
 f_{deceptive4} & 0.2 &  22  & 5497   & 63 & 7135  & 87   & 11257           \\ \hline 
                & 0.4 &  17  & 5839   & 62 & 8129 &  84  &  12516           \\ \hline
                & 0.6 &  6   & 6059  &  61 & 8377 &  75  &  14782           \\ \hline
                & 0.8 &  3   & 6059  &  55 & 8821 &  51  &  17983        \\ \hline
                & 1.0 &  2   & 3845  &  40 & 9490 &  3  &    6525        \\
                \hline \hline

 f_{3deceptive} & 0.2 &  24   & 10657 & 62 & 14454   & 98    &  17677      \\ \hline
                & 0.4 &  17   & 11283 & 63 & 14780   & 99    &  16964      \\ \hline
                & 0.6 &   6   & 10490 & 62 & 15389   & 90    &  21490      \\ \hline
                & 0.8 &   1   & 7993  & 47 & 16112   & 65    &  28418      \\ \hline
                & 1.0 &   1   & 17983 & 45 & 16694   & 9     &  27529     \\ \hline \hline
 F_{IsoP}       & 0.2 &   28     & 5888 &  15 & 6861   & 30    & 6894       \\ \hline 
                & 0.4 &   16     & 5808 & 17 &  7640  &  23   &  7949      \\ \hline
                & 0.6 &   21     & 5852 & 15 &  6661  &  17   &  6760      \\ \hline
                & 0.8 &   24     & 5828 & 19 & 6994   &  16   &  6744      \\ \hline
                & 1.0 &   21     & 5614 & 8  & 7243 &    13 &    5226    \\ \hline \hline
 F_{IsoTorus}   & 0.2 &  93    & 5436  & 100 & 6134   & 97    & 5706       \\ \hline 
                & 0.4 &  95      &5511  & 95 & 6268   & 95    & 5585       \\ \hline
                & 0.6 &  93      &5404  & 99 & 6348   & 94    & 5538       \\ \hline
                & 0.8 &  92      &5441  &100  & 6314  & 87    & 5627       \\ \hline
                & 1.0 &  89      &5456  & 98 &  6015  & 90    & 5507       \\ \hline 


\end{array}$
   \label{MT6FUNCS1}  
  \end{center}
 \end{table}

We evaluate the validity of our proposal in the next experiment. For the $5$
functions considered before and different values of the number of trees we run the
MT-FDA and determine the value $\mu$ in $[0.2,0.4,0.6,0.8,0,1.0]$ for which best
results are achieved. When the same successful rate is obtained  for more than one
value of $\mu$ we choose the smallest $\mu$ . Table~\ref{MT2_4_12FUNCS} presents the results
that confirm that our choice is correct even if there exist exceptions. In most of
the cases the $0.2$ value was the best or among the best assignments for $\mu$. This
experiment also allows to evaluate the influence of the number of trees in the
performance of the MT-FDA for the different functions.



\begin{table}[htbp ]
  \begin{center} 
   \caption{Results of the MT-FDA with different number of  components in the optimization of
    different functions  $n=36$, $T=0.15$ when the number of trees is increased}
   $\begin{array}{|r|r|r|r|r|r|r|r|r|r|r|}
\hline
  Function & N.Trees  &     \multicolumn{3}{|c|}{No Prior}  &  \multicolumn{3}{|c|}{Best Prior} &
      \multicolumn{3}{|c|}{Adap. Prior}  \\ \hline        
        &    & Best \mu     & succ.   & eval. & Best \mu & succ.  & eval. & Best \mu & succ.  & eval. \\ \hline \hline
 f_{deceptive3}& 2  & 0.4 &  100   & 3859  & 0.2  & 100 & 4440  & 0.2  &100  & 3851  \\ \hline
               & 4  & 0.2 &  98    & 4309  & 0.2  & 100   & 4810 &  0.2 & 100 & 4705      \\ \hline
               &12  & 0.2 &  66    & 5222  & 0.2  & 100   & 5858 &  0.2 & 100 & 6753 \\ \hline\hline
 f_{deceptive4} & 2  & 0.4 & 47 & 4760 & 0.8  & 71  & 6164 & 0.2 & 67 & 5916 \\ \hline
               & 4  & 0.2 &  34 &  5490 & 0.2  & 70 & 7400 & 0.4 & 72 & 8185     \\ \hline
               & 12 & 0.2 &  15 &  6059 & 0.2  & 75 & 8035 & 0.2 & 88 & 12916   \\   \hline \hline
 F_{IsoTorus}  & 2  & 0.2 & 94     & 5549 & 1.0 & 98 & 6515  & 0.2 & 92  &   5580    \\ \hline
               & 4  & 0.4 & 97     & 5490 & 0.6 & 100& 6345  & 0.4 & 98  &   5710  \\ \hline
               & 12 & 0.2 & 93     & 5608 & 0.4 & 99 & 6358  & 0.2 & 99  &   6056\\ \hline \hline
 F_{IsoP}      & 2  & 0.2 & 77     & 5100 & 0.2  &  41  &5654 & 1.0 &  70 & 5053      \\ \hline
               & 4 & 0.2 &  45     & 5573 & 0.4  &  32  &6089 & 0.6 &  38 & 5653   \\ \hline   
               & 12 & 0.2 & 23     & 6386 & 0.2 & 21 & 8183   & 0.2  & 18 & 8992     \\ \hline \hline
 f_{deceptive3}& 2  & 0.2 & 37     & 9370 & 0.2 & 66 & 12034 & 1.0& 63  & 11053       \\ \hline
               & 4 & 0.2 &  23     & 10338 & 0.2  & 66   & 13835     & 0.2 & 93 & 15942       \\ \hline   
               & 12 & 0.2 & 11     & 11172  &0.2  & 79   & 13911    &  0.2 & 95 & 20244    \\ \hline 
\end{array}$
    \label{MT2_4_12FUNCS}  
  \end{center}
 \end{table}
 

When priors are not used the algorithm is very sensitive to an increment in the
number of trees. If priors are incorporated the succesful rate is the same, or can
even increase when the number of trees is augmented, but at the expense of a higher number or function evaluations. The
exception is function $F_{IsoP}$, for this function results remarkably deteriorate
with the increment in the number of trees. $F_{IsoP}$ is an example of a function for
which a simple model gives better results. Also in the case of
$f_{deceptive3}$ it can be appreciated that a simple model can find the optimum with
a lesser number of evaluations. In the case of the other three functions results
improve when mixtures are used. 

\section{Conclusions and Future Research} \label{CONCLUSIONS}

 In this paper we have introduced a mutation like effect in the MT-FDA by
 using priors. We have proposed a way for doing adaptive mutation by relating
 the mutation rate of the FDA to the quality of the approximation and the strength of
 the components of the mixture model. This is an example of the number of
 interesting possibilities that arise in the integration of the theory of
 Graphical Models and Evolutionary Computation.  The results of the experiments show that the changes introduced to the MT-FDA
 improve its performance.

 As trends of future research we identify the creation of learning algorithms
 for the mixtures of trees that allow to determine the number of trees
 incrementally during the learning phase. Such types of algorithms would be
 able to adjust the number of components of the mixture to the characteristics
 of the data.  We plan to evaluate the convenience of using other types of
 priors, like structural priors that allow to change the structure of the
 trees, and not only the parameters, and "smoothing
with the marginal". This method is thought to allow the appearance of events with zero probability. The
probability distributions obtained from the data are smoothed with some more general
probability distribution by interpolating both distributions. This technique
 can be seen as a Dirichlet prior derived from the pairwise marginal distributions for the data set
\cite{Meila_and_Jordan:2000}. 

\bibliographystyle{abbrv}   
\bibliography{Thesbib}   
  
   
\end{document}   



